\chapter{Hyperfine and EFG parameters (menu 40)}
\label{ch:menu40}
\menulabel{40}

\section{Purpose}

The \code{\%eprnmr} module (manual section 7.51.3) prints the electric and
magnetic hyperfine structure: the A tensor per nucleus (Fermi contact,
spin dipole, second-order SOC), the electric field gradient with its
electron/nuclear decomposition, and the density at the nucleus.  This menu
reads those blocks and derives the NQR/Mossbauer observables: the nuclear
quadrupole coupling constant and the first-order quadrupole splitting from
a caller-supplied nuclear quadrupole moment.

\section{What you need}

An ORCA output with \code{\%eprnmr} and a \code{Nuclei} list; the
coordinates block must precede \code{\%eprnmr} when \code{Nuclei} is given
(measured -- otherwise the engine stops with ``nuclear properties are
requested but no coordinates have been read'').  The A tensor needs the
per-nucleus parameters (I, P via \code{PPP}/\code{III}); without them the
engine prints zeros.  Optionally answer the nuclear quadrupole moment
$Q$ (barn) for the quadrupole conversion.

\section{Walkthrough}

The example reads the all-electron CeF$_3$ probe (EFG principal values
0.6771983 / 0.7337722 / $-1.4109705$~a.u., $\eta$ = 0.0401,
Rho(0) = 727217~a.u.$^{-3}$, the A matrix zero because no nuclear
parameters were given) and converts with $Q$ = 0.5~barn:
$eQV_{zz}/h$ = $-165.7642$~MHz, $\Delta E_Q$ = 82.9043~MHz -- the
conversion agrees with the engine's own quadrupole block to $\sim10^{-6}$.

\lstinputlisting[style=fbkcode, firstline=53, lastline=69,
  title={The menu-40 example (the EFG principal values, the
  electron/nuclear decomposition and the quadrupole conversion)}]
  {examples/expected/m40_hyperfine.txt}

\section{What you get}

The report \code{<output>.hyperfine.fbk.md}: per nucleus the printed
nuclear parameters and the A tensor (A(iso), A(Tot), the FC/SD split);
the EFG principal values with $|V_{zz}|$, $\eta = (V_{xx}-V_{yy})/V_{zz}$
(the standard NQR ordering), the orientation and the V(El)/V(Nuc) sum
cross-check; Rho(0); and -- with $Q$ -- $eQV_{zz}/h$ and the
first-order $\Delta E_Q$.

\section{Conventions, cross-checks and boundaries (measured)}

\begin{readbox}
\begin{itemize}
  \item \textbf{The CASSCF route keeps the EFG but switches the A
    components off} (measured, method ``CASSCF/ALL STATES AVERAGE''):
    pair the DFT route for A with the CASSCF route for correlation-aware
    EFG/Rho(0) as needed.
  \item \textbf{Rho(0) is basis-domain dependent}: the all-electron probe
    reports 727216.9~a.u.$^{-3}$, the small-core ECP probe 0.075618828 on
    the same geometry -- never compare across domains (an isomer-shift
    reference must come from the same domain).
  \item \textbf{The EFG is sensitive to SCF micro-solutions} (measured:
    two identical-input runs reproduced digit-for-digit; a third differing
    by $2\times10^{-6}$~Eh moved $V_{\mathrm{Tot}}$ by $\sim10^{-4}$ relative)
    -- converged energy does not imply a reproducible EFG.
  \item \textbf{Cross-checks in the report}: V(El) + V(Nuc) reproduces
    V(Tot) to the printed digits, and the nuclear triplet V(Nuc) is
    identical across routes of the same geometry (both measured).
  \item \textbf{Boundary (termination)}: core-level photoemission (XPS)
    has no ORCA module (binary search: zero hits); the XES/SOC-XES strings
    in the relativistic CASSCF code belong to the ROCIS family (chapter
    46's domain) -- registered as a follow-up.
\end{itemize}
\end{readbox}
